% Imporditud: tools/import_eksperimendid.py
% Allikas: Eesti füüsikaolümpiaadi eksperimendi ülesannete
%   kogu (Rael Kalda, 2023), ülesanne Ü80, lahendus L80.
\setAuthor{EFO žürii}
\setRound{lõppvoor}
\setYear{1999}
\setNumber{G E2}
\setDifficulty{3}
\setTopic{elektriahelad}
\setVahendid{Kaks testrit koos juhtmetega, taskulambipatarei 4,5 V}
\setPunktid{12}
\setAllikas{Eksperimendi kogu Ü80, lahendus lk 68}
\setLahendusseis{toores}

\prob{Voltmeetri ja ampermeetri takistus}
Kasutades üht testrit voltmeetrina ja teist ampermeetrina, leida kummagi mõõteriista takistus sellisel mõõtepiirkonnal, mis on kõige sobivam kasutamiseks antud vooluallika korral. Testrit oommeetrina mitte kasutada.
\textit{Märkus:} Valem testri mõõtemääramatuse arvutamiseks pinge mõõtmisel on 0,005x+ 2y ning voolutugevuse mõõtmisel 0,02x + 2y, kus x on mõõtmistulemus ja y --- vähim võimalik mõõtmistulemus antud mõõtepiirkonnal. Näiteks kui x = 4,56 V, siis mõõtemääramatus on (0,005 $\cdot$ 4,56 + 0,02) V.

\hint

\solu
G $E_2$

Kui ühendame ampermeetri ja voltmeetri jadamisi, siis mõõdab ampermeeter voltmeetrit läbiva voolu tugevust ning voltmeeter mõõdab pinget iseenesel. Pinge ja voolu suhe on voltmeetri takistus (2 p). Sobilik mõõtepiirkond voltmeetril on 20 V ja ampermeetril 200 $\mu$A (1 p). Kui ühendame ampermeetri ja voltmeetri rööbiti, siis mõõdab ampermeeter vooluallika lühisvoolu tugevust (voltmeetri suure takistuse tõttu on teda läbiv vool tühine) ning voltmeeter mõõdab pinget rööpühenduse otstel. Pinge ja voolu suhe on ampermeetri takistus (2 p). Mõõtepiirkondade valik: ampermeetril 10 A, voltmeetril 20 V või 2000 mV (1 p). Korrektsed mõõtmistulemused takistuse leidmiseks: voltmeetril 1 p, ampermeetril 1 p; korrektsed arvutused: kumbki 1 p, mõõtmistäpsuse hinnang --- kumbki 1 p.
\probend
